\chapter{Effets électroniques et intermédiaires réactifs}\label{ch:b1:electronic-effects}

L'acide trifluoroéthanoïque, \ce{CF3COOH}, est un acide environ dix
mille fois plus fort que l'acide éthanoïque, \ce{CH3COOH}, bien que la
liaison O--H acide soit à trois liaisons des atomes de fluor. Le phénol
est un million de fois plus acide que l'éthanol, et une molécule
d'ammoniac liée à un cycle benzénique --- l'aniline --- est un million de
fois moins basique qu'une molécule liée à un groupe méthyle. La chimie
organique abonde en contrastes de ce genre, et deux effets expliquent la
plupart d'entre eux~: l'attraction exercée par les atomes
électronégatifs le long des liaisons $\sigma$, et la délocalisation des
doublets d'électrons sur les systèmes $\pi$. Ces deux mêmes effets
décident quelles liaisons se rompent, quels intermédiaires se forment et
où attaquent les réactifs~: ils sont la grammaire des \omterm{def:b1:elementary-steps:elementary-step}{mécanismes
réactionnels} des chapitres suivants.

\begin{recall}
L'\omterm{def:b1:periodicity:electronegativity}{électronégativité} et ses évolutions (\cref{ch:b1:periodicity})~; les
\omterm{def:b1:lewis-resonance-vsepr:lewis}{schémas de Lewis}, les \omterm{def:b1:lewis-resonance-vsepr:formal-charge}{charges formelles} et la mésomérie
(\cref{ch:b1:lewis-resonance-vsepr})~; le $\mathrm{p}K_a$ et la force
des acides et des bases (\cref{ch:b1:predominant-reaction}). Le volume
scolaire (année 12) a dessiné des \omterm{def:b1:electronic-effects:curly-arrow}{flèches courbes} allant d'un \omterm{def:b1:electronic-effects:nucleophile}{nucléophile}
à un \omterm{def:b1:electronic-effects:nucleophile}{électrophile}.
\end{recall}

\section{Effets inductifs}

\begin{definition}[Effet inductif]\label{def:b1:electronic-effects:inductive}
L'\emph{effet inductif}\index{effet inductif} est le déplacement de
densité électronique le long des liaisons $\sigma$ vers un atome ou un
groupe électronégatif. Un \omterm{def:b1:electronic-effects:mesomeric}{groupe attracteur} (F, Cl, O, \ce{NO2},
\ce{CF3}) exerce un effet $-I$~; les groupes alkyle et les atomes
chargés négativement repoussent les électrons, un effet $+I$. L'effet
s'affaiblit rapidement avec le nombre de liaisons intermédiaires.
\end{definition}

\begin{proposition}[Effet inductif et acidité]\label{prop:b1:electronic-effects:inductive-pka}
Un \omterm{def:b1:electronic-effects:mesomeric}{groupe attracteur} proche du site acide d'un acide carboxylique
stabilise l'ion carboxylate en étalant sa charge négative, et abaisse le
$\mathrm{p}K_a$~: plus les \omterm{def:b1:electronic-effects:mesomeric}{groupes attracteurs} sont nombreux et proches,
plus l'acide est fort.
\end{proposition}

\begin{proof}
$\mathrm{p}K_a$ mesure la position de $\mathrm{RCOOH} \rightleftharpoons
\mathrm{RCOO^-} + \ce{H+}$. Un groupe qui attire la densité électronique
hors du carboxylate étale sa charge sur davantage d'atomes, ce qui
abaisse son énergie plus que celle de l'acide neutre~; l'équilibre se
déplace vers la droite et $K_a$ augmente. La série mesurée le confirme.
\end{proof}

\begin{omfigure}
\begin{tikzpicture}
\begin{axis}[
  xbar, width=0.78\linewidth, height=5.2cm, xmin=0, xmax=5.4,
  symbolic y coords={TFA,TCA,DCA,MCA,AA}, ytick=data,
  yticklabels={\ce{CF3COOH}, \ce{CCl3COOH}, \ce{CHCl2COOH}, \ce{CH2ClCOOH}, \ce{CH3COOH}},
  xlabel={$\mathrm{p}K_a$ à \qty{25}{\celsius}}, bar width=9pt,
  nodes near coords, nodes near coords align={horizontal},
  every node near coord/.append style={font=\scriptsize},
  tick label style={font=\small}, label style={font=\small}, enlarge y limits=0.15]
\addplot[fill=omDef!45, draw=omDef] coordinates {(0.52,TFA) (0.51,TCA) (1.35,DCA) (2.87,MCA) (4.76,AA)};
\end{axis}
\end{tikzpicture}

{\small Les acides éthanoïques halogénés. Chaque chlore abaisse le
$\mathrm{p}K_a$~; avec trois chlores, l'acide est environ dix mille fois
plus fort que l'acide éthanoïque. Les acides trifluoroéthanoïque et
trichloroéthanoïque sont aussi forts l'un que l'autre, compte tenu de
l'incertitude des mesures sur des acides aussi forts.}
% ledger: pka:ethanoic, pka:chloroethanoic, pka:dichloroethanoic, pka:trichloroethanoic, pka:trifluoroethanoic
\end{omfigure}

Les groupes alkyle agissent en sens inverse, un peu~: l'acide
propanoïque ($\mathrm{p}K_a
= 4.89$) est légèrement plus faible que l'acide éthanoïque (4.76), et
l'acide méthanoïque (3.74), dépourvu de groupe alkyle, est plus fort que
les deux.
% ledger: pka:propanoic, pka:methanoic

\section{Effets mésomères}

Lorsqu'un atome porteur d'un \omterm{def:b1:lewis-resonance-vsepr:lewis}{doublet non liant}, ou d'une charge, est lié
à un système $\pi$, ses électrons peuvent être partagés avec ce
système~; les \omterm{def:b1:lewis-resonance-vsepr:resonance}{formules mésomères} (\cref{ch:b1:lewis-resonance-vsepr})
décrivent ce partage.

\begin{definition}[Effet mésomère, groupes donneurs et attracteurs]\label{def:b1:electronic-effects:mesomeric}
L'\emph{effet mésomère}\index{effet mesomere@effet mésomère} est le déplacement de densité
électronique à travers un système $\pi$ conjugué, décrit par des
\omterm{def:b1:lewis-resonance-vsepr:resonance}{formules mésomères}. Un \emph{groupe donneur}\index{groupe donneur} cède
un \omterm{def:b1:lewis-resonance-vsepr:lewis}{doublet non liant} au système, un effet $+M$ (\ce{-OH}, \ce{-O^-},
\ce{-NH2}, \ce{-OR})~; un \emph{groupe attracteur}\index{groupe attracteur}
retire des électrons au système par une liaison $\pi$ vers un atome
électronégatif, un effet $-M$ (\ce{-NO2},
\ce{-C=O}, \ce{-C#N}).
\end{definition}

\begin{omfigure}
\setchemfig{atom sep=2em}
\schemestart
\chemfig{CH_3-C(=[1]O)-[7]O^{-}}
\arrow{<->}
\chemfig{CH_3-C(-[1]O^{-})=[7]O}
\schemestop
\hspace{2.5em}
\schemestart
\chemfig{*6(=-=(-O^{-})-=-)}
\arrow{<->}
\chemfig{*6((=O)-=-\chemabove{}{\scriptstyle\ominus}-=-)}
\schemestop

\medskip
\schemestart
\chemfig{H_3C-C(=[1]O)-[7]NH_2}
\arrow{<->}
\chemfig{H_3C-C(-[1]O^{-})=[7]NH_2^{+}}
\schemestop

{\small \omterm{def:b1:lewis-resonance-vsepr:resonance}{Formules mésomères} de l'ion éthanoate (la charge partagée par
deux oxygènes équivalents), de l'ion phénolate (la charge étalée dans le
cycle~; une seule des trois formules du cycle est représentée) et d'un
amide (le doublet de l'azote partagé avec le C=O, ce qui fait des amides
de mauvaises bases).}
\end{omfigure}

\begin{proposition}[Effet mésomère et acidité]\label{prop:b1:electronic-effects:mesomeric-pka}
Un acide dont la base conjuguée étale sa charge par mésomérie est plus
fort qu'un acide dont la base ne le peut pas~: les acides carboxyliques
($\mathrm{p}K_a$ voisin de 4--5) sont bien plus forts que les alcools
(environ 16), et le phénol (10.0) se situe entre les deux. Un \omterm{def:b1:electronic-effects:mesomeric}{groupe
attracteur} en position \textit{ortho} ou \textit{para} par rapport au OH
d'un phénol le renforce davantage qu'en position \textit{méta}.
\end{proposition}

\begin{proof}
Dans un alcoolate, la charge est portée par un seul oxygène~; dans un
carboxylate, elle est partagée par deux oxygènes~; dans un phénolate,
elle est partagée par l'oxygène et trois carbones du cycle, moins
électronégatifs que l'oxygène, d'où un gain moindre. Un groupe
\ce{NO2} sur le carbone en \textit{ortho} ou en \textit{para} peut
prendre lui-même la charge grâce à une \omterm{def:b1:lewis-resonance-vsepr:resonance}{formule mésomère} comportant
\ce{C=N+(-O^-)-O^-}~; en
\textit{méta}, aucune \omterm{def:b1:lewis-resonance-vsepr:resonance}{formule mésomère} n'amène la charge sur le carbone
qui porte \ce{NO2}, et seul l'\omterm{def:b1:electronic-effects:inductive}{effet inductif} subsiste. Les valeurs
mesurées~: 2-nitrophénol 7.23, 4-nitrophénol 7.16, 3-nitrophénol 8.36,
contre 9.99 pour le phénol et 15.9 pour l'éthanol.
\end{proof}
% ledger: pka:ethanol, pka:phenol, pka:2-nitrophenol, pka:3-nitrophenol, pka:4-nitrophenol

Le même raisonnement s'applique aux bases. La méthylamine
($\mathrm{p}K_a$ de son ion ammonium 10.66) est une base plus forte que
l'ammoniac (9.25, \cref{ch:b1:predominant-reaction}),
le groupe méthyle repoussant les électrons vers l'azote~; dans
l'aniline, le \omterm{def:b1:lewis-resonance-vsepr:lewis}{doublet non liant} est partagé avec le cycle et l'ion
ammonium est bien plus acide (4.60). La pyridine (5.23) garde son
doublet, mais sur un azote engagé dans un cycle plan, de caractère
\textit{s} plus marqué, donc moins disponible que dans une amine.
% ledger: pka:methylammonium, pka:anilinium, pka:pyridinium, dfg:NH4+_ao

\section{Rompre une liaison~: les intermédiaires réactifs}

\begin{definition}[Homolyse et hétérolyse]\label{def:b1:electronic-effects:cleavage}
Une \omterm{def:b1:lewis-resonance-vsepr:lewis}{liaison covalente} se rompt par \emph{homolyse}\index{homolyse}
lorsque chaque atome garde l'un des deux électrons liants, ce qui donne
deux radicaux, et par \emph{hétérolyse}\index{heterolyse@hétérolyse} lorsqu'un atome les garde
tous les deux, ce qui donne un cation et un anion.
\end{definition}

\begin{definition}[Intermédiaires réactifs]\label{def:b1:electronic-effects:intermediates}
Un \emph{intermédiaire réactif}\index{intermediaire reactif@intermédiaire réactif} est une espèce formée dans
une étape d'un mécanisme et consommée dans une étape ultérieure, trop
réactive pour s'accumuler (\cref{ch:b1:elementary-steps}). Un
\emph{carbocation}\index{carbocation} possède un carbone portant trois
liaisons et une charge positive (six \omterm{def:b1:quantum-numbers:configuration}{électrons de valence})~; un
\emph{carbanion}\index{carbanion}, un carbone portant trois liaisons, un
\omterm{def:b1:lewis-resonance-vsepr:lewis}{doublet non liant} et une charge négative~; un \emph{radical}\index{radical},
un atome porteur d'un \omterm{def:b1:quantum-numbers:unpaired}{électron célibataire}.
\end{definition}

\begin{definition}[Classe d'un carbone]\label{def:b1:electronic-effects:carbon-class}
Un carbone lié à un seul autre carbone est un \emph{carbone
primaire}\index{carbone primaire}, à deux un \emph{carbone
secondaire}\index{carbone secondaire}, à trois un \emph{carbone
tertiaire}\index{carbone tertiaire}. Un \omterm{def:b1:electronic-effects:intermediates}{carbocation}, un \omterm{def:b1:electronic-effects:intermediates}{radical} ou un
halogénoalcane prend la classe du carbone concerné.
\end{definition}

\begin{omfigure}
\begin{tikzpicture}[font=\small]
  % carbocation: planar, empty p orbital
  \begin{scope}
    \fill[omProp!15, draw=omProp] (0,0.15) ellipse (0.18 and 0.75);
    \fill[omProp!15, draw=omProp] (0,-0.15) ellipse (0.18 and 0.75);
    \draw[thick] (0,0) -- (-1.0,-0.25) node[left] {R};
    \draw[thick] (0,0) -- (0.95,-0.35) node[right] {R};
    \draw[thick] (0,0) -- (0.35,0.45) node[above right] {R};
    \node[fill=white, inner sep=1pt] at (0,0) {\ce{C+}};
    \node at (0,-1.75) {carbocation~: plan,};
    \node at (0,-2.15) {orbitale \textit{p} vide};
  \end{scope}
  % radical: nearly planar, one electron
  \begin{scope}[xshift=4.6cm]
    \draw[omMeth] (0,0.15) ellipse (0.18 and 0.75);
    \draw[omMeth] (0,-0.15) ellipse (0.18 and 0.75);
    \fill (0,0.55) circle (0.05);
    \draw[thick] (0,0) -- (-1.0,-0.25) node[left] {R};
    \draw[thick] (0,0) -- (0.95,-0.35) node[right] {R};
    \draw[thick] (0,0) -- (0.35,0.45) node[above right] {R};
    \node[fill=white, inner sep=1pt] at (0,0) {C};
    \node at (0,-1.75) {radical~: presque plan,};
    \node at (0,-2.15) {un électron};
  \end{scope}
  % carbanion: pyramidal with a lone pair
  \begin{scope}[xshift=9.2cm]
    \draw[thick] (0,0) -- (-0.95,-0.55) node[left] {R};
    \draw[thick] (0,0) -- (0.95,-0.55) node[right] {R};
    \draw[thick] (0,0) -- (0.2,-0.9) node[below] {R};
    \fill[omDef!15, draw=omDef] (0,0.55) ellipse (0.22 and 0.45);
    \fill (-0.07,0.65) circle (0.045); \fill (0.07,0.65) circle (0.045);
    \node[fill=white, inner sep=1pt] at (0,0) {\ce{C-}};
    \node at (0,-1.75) {carbanion~: pyramidal,};
    \node at (0,-2.15) {doublet non liant};
  \end{scope}
\end{tikzpicture}

{\small Géométries des trois \omterm{def:b1:electronic-effects:intermediates}{intermédiaires réactifs} du carbone (R~: H ou
un groupe alkyle). Le \omterm{def:b1:electronic-effects:intermediates}{carbocation}, plan, peut être attaqué sur l'une ou
l'autre face, point qui compte pour la stéréochimie du
\cref{ch:b1:nucleophilic-substitution}.}
\end{omfigure}

\begin{proposition}[Stabilité des carbocations]\label{prop:b1:electronic-effects:carbocation-order}
Les \omterm{def:b1:electronic-effects:intermediates}{carbocations} sont stabilisés par les groupes qui leur cèdent des
électrons~: tertiaire $>$ secondaire $>$ primaire $>$ méthyle~; un
\omterm{def:b1:electronic-effects:intermediates}{carbocation} voisin d'une liaison double (allylique) ou d'un cycle
benzénique (benzylique) est en outre stabilisé par mésomérie, et un
\omterm{def:b1:electronic-effects:intermediates}{carbocation} voisin d'un atome porteur d'un \omterm{def:b1:lewis-resonance-vsepr:lewis}{doublet non liant}
(\ce{R-O-CH2+}) l'est encore davantage.
\end{proposition}

\begin{proof}
Chaque groupe alkyle repousse de la densité électronique vers le carbone
pauvre en électrons (effet $+I$, et une interaction apparentée de ses
liaisons C--H avec l'orbitale vide, décrite dans le volume de deuxième
année). Un cation allylique possède deux \omterm{def:b1:lewis-resonance-vsepr:resonance}{formules mésomères} qui placent
la charge sur l'un ou l'autre carbone terminal~; un cation benzylique en
possède quatre, dont trois dans le cycle~; un \omterm{def:b1:lewis-resonance-vsepr:lewis}{doublet non liant} de
l'oxygène forme une quatrième liaison avec le carbone cationique,
\ce{R-O+=CH2}, où chaque atome a son octet.
\end{proof}

Les \omterm{def:b1:electronic-effects:intermediates}{carbanions} suivent l'ordre inverse (méthyle $>$ primaire $>$
secondaire
$>$ tertiaire) et sont stabilisés par les \omterm{def:b1:electronic-effects:mesomeric}{groupes attracteurs}~; les
radicaux suivent le même ordre que les \omterm{def:b1:electronic-effects:intermediates}{carbocations}, de façon moins
marquée.

\section{Flèches courbes, nucléophiles et électrophiles}

\begin{definition}[Flèches courbes]\label{def:b1:electronic-effects:curly-arrow}
Dans un mécanisme, une \emph{flèche courbe}\index{fleche courbe@flèche courbe} représente le
mouvement d'un doublet d'électrons~: elle part d'un \omterm{def:b1:lewis-resonance-vsepr:lewis}{doublet non liant} ou
d'une liaison et aboutit sur un atome (où se forme un \omterm{def:b1:lewis-resonance-vsepr:lewis}{doublet non liant}
ou une nouvelle liaison) ou entre deux atomes (une nouvelle liaison).
Une \emph{demi-flèche}\index{demi-fleche@demi-flèche} représente le mouvement d'un seul électron,
dans les réactions radicalaires.
\end{definition}

\begin{method}[Dessiner des flèches courbes]\label{met:b1:electronic-effects:curly-arrows}
\begin{enumerate}
  \item Partir des électrons~: un \omterm{def:b1:lewis-resonance-vsepr:lewis}{doublet non liant}, une liaison $\pi$ ou
        une liaison $\sigma$, jamais d'une charge positive ni d'un atome
        dépourvu d'électrons.
  \item Aboutir là où va le doublet~: une nouvelle liaison entre les deux
        atomes, ou un \omterm{def:b1:lewis-resonance-vsepr:lewis}{doublet non liant} sur l'atome le plus
        électronégatif lorsqu'une liaison se rompt.
  \item Compter les charges après chaque étape~: l'atome qui cède un
        \omterm{def:b1:lewis-resonance-vsepr:lewis}{doublet non liant} à une nouvelle liaison gagne $+1$, l'atome qui
        reçoit un doublet sous forme de \omterm{def:b1:lewis-resonance-vsepr:lewis}{doublet non liant} gagne $-1$~;
        la charge totale se conserve.
  \item Ne jamais dépasser l'octet sur C, N, O~: quand un doublet arrive,
        un autre doit partir.
\end{enumerate}
\end{method}

\begin{omfigure}
\setchemfig{atom sep=2.2em}
\schemestart
\chemfig{H_3C-C(-[2]H)(-[6]H)-[@{b}]@{br}Br}
\arrow{->[hétérolyse]}[,1.5]
\chemfig{H_3C-C(-[2]H)(-[6]H)^{+}}\+\chemfig{Br^{-}}
\schemestop
\chemmove{\draw[->, shorten <=2pt, shorten >=2pt](b) .. controls +(90:6mm) and +(90:6mm) .. (br);}

\medskip
\schemestart
\chemfig{H@{o}O^{-}}\+\chemfig{@{h}H-@{oc}O-H}
\arrow{<=>}
\chemfig{H_2O}\+\chemfig{HO^{-}}
\schemestop
\chemmove{\draw[->, thick, shorten <=2pt, shorten >=2pt](o) .. controls +(-80:8mm) and +(-100:8mm) .. (h);}
\par\vspace{5mm}

{\small \omterm{def:b1:electronic-effects:curly-arrow}{Flèches courbes}. En haut~: \omterm{def:b1:electronic-effects:cleavage}{hétérolyse} de la liaison C--Br, le
doublet allant au brome. En bas~: transfert d'un proton de l'eau à l'ion
hydroxyde, écrit comme l'attaque d'un \omterm{def:b1:lewis-resonance-vsepr:lewis}{doublet non liant} de la base sur
l'hydrogène (la liaison O--H de l'eau se rompt alors, son doublet
restant sur l'oxygène~; la seconde flèche est laissée à
l'\cref{exo:b1:electronic-effects:4}).}
\end{omfigure}

\begin{definition}[Nucléophile, électrophile, groupe partant]\label{def:b1:electronic-effects:nucleophile}
Un \emph{nucléophile}\index{nucleophile@nucléophile} est une espèce qui cède un doublet
d'électrons à un atome autre que l'hydrogène pour former une liaison
(une \omterm{def:b1:lewis-resonance-vsepr:lewis-acid}{base de Lewis})~; un \emph{électrophile}\index{electrophile@électrophile} est une espèce qui
accepte un tel doublet (un \omterm{def:b1:lewis-resonance-vsepr:lewis-acid}{acide de Lewis}, souvent un carbone porteur
d'une charge partielle positive). La \emph{nucléophilie}\index{nucleophilie@nucléophilie} d'une
espèce est sa réactivité en tant que nucléophile, mesurée par des
\omterm{def:b1:rate-laws:order}{constantes de vitesse}. Un \emph{groupe partant}\index{groupe partant} est
le groupe qui s'en va avec le \omterm{def:b1:lewis-resonance-vsepr:lewis}{doublet liant} lorsqu'une liaison au
carbone se rompt de façon hétérolytique.
\end{definition}

\begin{proposition}[Nucléophilie et basicité]\label{prop:b1:electronic-effects:nucleophilicity-basicity}
Pour des \omterm{def:b1:electronic-effects:nucleophile}{nucléophiles} qui attaquent par le même atome, la \omterm{def:b1:electronic-effects:nucleophile}{nucléophilie}
suit la basicité~: \ce{CH3O-} $>$ \ce{OH-} $>$ \ce{CH3COO-} $>$
\ce{H2O}. Ce parallèle est en défaut le long d'une colonne du tableau
périodique dans les \omterm{def:b1:intermolecular-forces-solvents:solvent-classes}{solvants protiques}, où les \omterm{def:b1:electronic-effects:nucleophile}{nucléophiles} gros,
polarisables et peu solvatés réagissent plus vite bien qu'ils soient des
bases plus faibles (\ce{I-} $>$ \ce{Br-} $>$ \ce{Cl-} $>$ \ce{F-}), et
pour les bases encombrées, qui sont de mauvais \omterm{def:b1:electronic-effects:nucleophile}{nucléophiles}.
\end{proposition}

\begin{proof}
La basicité mesure l'équilibre de l'attaque sur un proton, la
\omterm{def:b1:electronic-effects:nucleophile}{nucléophilie} la vitesse de l'attaque sur un carbone~; toutes deux
augmentent avec la disponibilité du doublet, d'où le parallèle pour un
même atome. La vitesse dépend aussi de l'énergie nécessaire pour
débarrasser le \omterm{def:b1:electronic-effects:nucleophile}{nucléophile} de son solvant et de la déformation de son
nuage électronique vers un carbone encombré~: les petits anions comme
\ce{F-} sont fortement liés aux molécules de \omterm{def:b1:intermolecular-forces-solvents:solvent-classes}{solvant protique}
(\cref{ch:b1:intermolecular-forces-solvents}), et les groupes
volumineux gênent l'approche d'un carbone, mais pas celle d'un proton.
\end{proof}

Un bon \omterm{def:b1:electronic-effects:nucleophile}{groupe partant} est une \omterm{def:b1:predominant-reaction:strong-acid}{base faible}, stable une fois qu'il emporte
le doublet~: \ce{I-}, \ce{Br-}, \ce{Cl-}, l'eau, les sulfonates~;
\ce{OH-} et \ce{RO-}, \omterm{def:b1:predominant-reaction:strong-acid}{bases fortes}, sont de mauvais \omterm{def:b1:electronic-effects:nucleophile}{groupes partants} ---
c'est pourquoi il faut activer les alcools avant une substitution
(\cref{ch:b1:alcohol-activation}).

\section{Exercices}

\begin{exercise}[$\star$]\label{exo:b1:electronic-effects:1}
Classer chaque carbone du 2-méthylbutane et du 2-bromo-2-méthylpropane
en primaire, secondaire, tertiaire ou quaternaire (lié à quatre
carbones). Quel \omterm{def:b1:electronic-effects:intermediates}{carbocation} se forme le plus facilement à partir de
chaque bromure de formule \ce{C4H9Br}~?
\end{exercise}

\begin{exercise}[$\star$]\label{exo:b1:electronic-effects:2}
À partir des valeurs de $\mathrm{p}K_a$, calculer le rapport des
\omterm{def:b1:predominant-reaction:acidity-constant}{constantes d'acidité} des acides chloroéthanoïque et éthanoïque, et des
acides trifluoroéthanoïque et éthanoïque.
\end{exercise}

\begin{exercise}[$\star$]\label{exo:b1:electronic-effects:3}
Dessiner les \omterm{def:b1:lewis-resonance-vsepr:resonance}{formules mésomères} de l'ion éthanoate, de l'ion
4-nitrophénolate (avec la charge sur un oxygène du groupe nitro) et du
cation allyle \ce{CH2=CH-CH2+}.
\end{exercise}

\begin{exercise}[$\star$]\label{exo:b1:electronic-effects:4}
Compléter les \omterm{def:b1:electronic-effects:curly-arrow}{flèches courbes} du transfert de proton de la figure et de
\ce{NH3 + H3O+ -> NH4+ + H2O}. Vérifier les charges.
\end{exercise}

\begin{exercise}[$\star\star$]\label{exo:b1:electronic-effects:5}
Classer par acidité croissante~: l'éthanol, le phénol, l'acide
éthanoïque, le 4-nitrophénol, l'acide trichloroéthanoïque, l'eau
($\mathrm{p}K_a$ du couple \ce{H2O}/\ce{OH-}, 14.0). Justifier chaque
étape par un argument inductif ou mésomère.
\end{exercise}

\begin{exercise}[$\star\star$]\label{exo:b1:electronic-effects:6}
Classer par basicité croissante~: l'aniline, la méthylamine,
l'ammoniac, la pyridine. Expliquer la position de l'aniline à l'aide
d'une \omterm{def:b1:lewis-resonance-vsepr:resonance}{formule mésomère}.
\end{exercise}

\begin{exercise}[$\star\star$]\label{exo:b1:electronic-effects:7}
Pourquoi le 3-nitrophénol est-il moins acide que le 4-nitrophénol, mais
plus acide que le phénol~? Dessiner les \omterm{def:b1:lewis-resonance-vsepr:resonance}{formules mésomères} qui
justifient les deux affirmations.
\end{exercise}

\begin{exercise}[$\star\star$]\label{exo:b1:electronic-effects:8}
Classer les \omterm{def:b1:electronic-effects:intermediates}{carbocations} \ce{CH3+}, \ce{(CH3)3C+}, \ce{CH2=CH-CH2+},
\ce{CH3CH2+}, \ce{C6H5CH2+}, \ce{CH3OCH2+}, et justifier.
\end{exercise}

\begin{exercise}[$\star\star$]\label{exo:b1:electronic-effects:9}
Identifier le \omterm{def:b1:electronic-effects:nucleophile}{nucléophile} et l'\omterm{def:b1:electronic-effects:nucleophile}{électrophile} dans~: \ce{CH3Br + OH-}~;
\ce{H2O + H+}~; l'éthanal avec un ion cyanure~; \ce{BF3 + NH3}.
Dessiner les \omterm{def:b1:electronic-effects:curly-arrow}{flèches courbes}.
\end{exercise}

\begin{exercise}[$\star\star\star$]\label{exo:b1:electronic-effects:10}
Une solution contient \qty{0.010}{mol/L} d'acide trichloroéthanoïque. La
formule de l'\omterm{def:b1:predominant-reaction:strong-acid}{acide faible}, $\mathrm{pH} = \frac12(\mathrm{p}K_a - \log C)$,
est-elle valable~? Calculer correctement le pH et le \omterm{def:b1:predominant-reaction:dissociation}{coefficient de
dissociation}.
\end{exercise}

\begin{exercise}[$\star\star\star$]\label{exo:b1:electronic-effects:11}
Les $\mathrm{p}K_a$ des alcools mesurés dans l'eau augmentent du
méthanol (15.5) à l'éthanol (15.9), au propan-2-ol (17.1) et au
2-méthylpropan-2-ol (19.2). Donner deux explications, l'une fondée sur
l'\omterm{def:b1:electronic-effects:inductive}{effet inductif} des groupes alkyle sur l'alcoolate, l'autre sur la
\omterm{def:b1:intermolecular-forces-solvents:dissolution}{solvatation} de l'alcoolate.
% ledger: pka:methanol, pka:ethanol, pka:2-propanol, pka:tBuOH
\end{exercise}

\begin{exercise}[$\star\star\star$]\label{exo:b1:electronic-effects:12}
Expliquer pourquoi les amides ne sont ni basiques sur l'azote ni
\omterm{def:b1:electronic-effects:nucleophile}{nucléophiles}, alors que les amines sont l'un et l'autre, à l'aide de la
\omterm{def:b1:lewis-resonance-vsepr:resonance}{formule mésomère} de l'amide. Quel atome d'un amide est protoné en milieu
\omterm{def:b1:predominant-reaction:strong-acid}{acide fort}~?
\end{exercise}

\section{Problème~: pourquoi l'acide trifluoroéthanoïque est si fort}

\begin{problem}[{Problème du week-end --- la série inductive des acides
halogénés, les carboxylates face aux alcoolates, les trois nitrophénols,
un calcul par la réaction prépondérante, et le rapport des constantes
d'acidité des acides trifluoroéthanoïque et éthanoïque}]
\label{pb:b1:electronic-effects:1}
Données ($\mathrm{p}K_a$ à \qty{25}{\celsius})~: acide éthanoïque 4.76,
chloroéthanoïque 2.87, dichloroéthanoïque 1.35, trichloroéthanoïque 0.51,
trifluoroéthanoïque 0.52, éthanol 15.9, phénol 9.99, 2-nitrophénol 7.23,
3-nitrophénol 8.36, 4-nitrophénol 7.16~; $\mathrm{p}K_e = 14.00$.

\medskip
\noindent\textbf{Partie I --- La série inductive.}
\begin{enumerate}
  \item Calculer le rapport des $K_a$ pour chaque étape, de l'acide
        éthanoïque à l'acide trichloroéthanoïque.
  \item L'effet de chaque chlore est-il le même~? Commenter.
  \item Expliquer l'évolution par l'\omterm{def:b1:electronic-effects:inductive}{effet inductif}.
  \item Pourquoi un chlore porté, sur une chaîne plus longue, par le
        carbone le plus éloigné de COOH n'aurait-il presque aucun effet~?
  \item Le fluor est plus électronégatif que le chlore. Que montre la
        comparaison des acides trifluoré et trichloré, et pourquoi est-il
        difficile de mesurer de tels $\mathrm{p}K_a$ avec précision~?
  \item Sous quelle forme se trouve chacun de ces acides à pH 3~?
\end{enumerate}

\medskip
\noindent\textbf{Partie II --- Carboxylates et alcoolates.}
\begin{enumerate}[resume]
  \item Dessiner les deux \omterm{def:b1:lewis-resonance-vsepr:resonance}{formules mésomères} de l'ion éthanoate.
  \item Que prévoient-elles pour les longueurs des deux liaisons C--O de
        l'ion~?
  \item Pourquoi l'ion éthanolate n'a-t-il pas de formules équivalentes~?
  \item Calculer le rapport des $K_a$ de l'acide éthanoïque et de
        l'éthanol.
  \item Placer le phénol entre les deux et justifier à l'aide des
        \omterm{def:b1:lewis-resonance-vsepr:resonance}{formules mésomères} du phénolate.
  \item Lesquels de l'éthanol, du phénol et de l'acide éthanoïque
        réagissent presque totalement avec une solution
        d'hydrogénocarbonate ($\mathrm{p}K_a$ 6.37 pour
        \ce{CO2,H2O}/\ce{HCO3-})~?
\end{enumerate}

\medskip
\noindent\textbf{Partie III --- Les nitrophénols.}
\begin{enumerate}[resume]
  \item Classer les trois nitrophénols et le phénol par acidité.
  \item Dessiner une \omterm{def:b1:lewis-resonance-vsepr:resonance}{formule mésomère} du 4-nitrophénolate portant la
        charge sur le groupe nitro.
  \item Montrer qu'une telle formule n'existe pas pour le
        3-nitrophénolate.
  \item Pourquoi le 3-nitrophénol est-il quand même plus acide que le
        phénol~?
  \item Le 2-nitrophénol est un peu moins acide que le 4-nitrophénol,
        bien que tous deux présentent l'\omterm{def:b1:electronic-effects:mesomeric}{effet mésomère}~; proposer une
        explication faisant intervenir son OH et le groupe nitro voisin.
  \item À pH 7.5, quels nitrophénols sont surtout sous forme anionique~?
        Pourquoi cela fait-il du 4-nitrophénol un \omterm{def:b1:titration-methods:indicator}{indicateur coloré}~?
\end{enumerate}

\medskip
\noindent\textbf{Partie IV --- L'acide trifluoroéthanoïque dans l'eau.}
\begin{enumerate}[resume]
  \item Calculer le pH de l'acide trifluoroéthanoïque à
        \qty{0.010}{mol/L} (sans supposer une faible dissociation).
  \item Calculer le pH de l'acide éthanoïque à \qty{0.010}{mol/L}.
  \item Écrire la réaction de l'acide trifluoroéthanoïque avec
        l'éthanoate de sodium et calculer sa constante.
  \item Donner le rapport $K_a(\ce{CF3COOH})/K_a(\ce{CH3COOH})$, sous
        forme de puissance de dix et sous forme de nombre.
\end{enumerate}
\end{problem}
